\documentclass{amsart}

\usepackage[T1]{fontenc}
\usepackage{lmodern}
\usepackage{microtype}
\usepackage{amsmath,amssymb}
\usepackage[hidelinks]{hyperref}

\hypersetup{
  pdftitle={A counterexample to dimension-two stabilization for Mn(S7)},
  pdfauthor={Owen Shen},
  pdfsubject={Additively idempotent semirings and generated varieties},
  pdfkeywords={additively idempotent semiring, matrix semiring, identity, variety}
}

\newcommand{\M}{\mathbf{M}}
\newcommand{\V}{\mathsf{V}}

\newtheorem{theorem}{Theorem}

\title[Dimension-two stabilization for $\M_n(S_7)$]
{A counterexample to dimension-two stabilization for $\M_n(S_7)$}

\author{Owen Shen}
\address{Operations Research Center, Massachusetts Institute of Technology,
Cambridge, Massachusetts 02139, USA}
\email{owenshen@mit.edu}
\date{August 2026}

\subjclass[2020]{Primary 16Y60; Secondary 03C05, 08B15}
\keywords{additively idempotent semiring, matrix semiring, identity, variety}

\begin{document}

\begin{abstract}
Jiao and Ren asked whether
$\V(\M_n(S_7))=\V(\M_{n+1}(S_7))$ for every $n\geq 2$ and suggested
that the chain may stabilize at dimension two. We show that this is false:
the identity
\[
 x+y+x^2+y^2 \approx x+y+x^2+y^2+yx
\]
holds in $\M_2(S_7)$ and fails in $\M_3(S_7)$. Consequently,
$\V(\M_2(S_7))\subsetneq\V(\M_3(S_7))$. The same identity fails in
$\M_n(S_7)$ for every $n\geq 3$.
\end{abstract}

\maketitle

An \emph{additively idempotent semiring} is an algebra $(S,+,\cdot)$
whose additive reduct is a commutative idempotent semigroup, whose
multiplicative reduct is a semigroup, and in which multiplication distributes
over addition. For such an algebra $A$, let $\V(A)$ denote the variety it
generates. Let $S_7=\{1,a,\infty\}$ be the additively idempotent semiring
with Cayley tables
\[
\begin{array}{c|ccc}
 + & 1 & a & \infty \\ \hline
 1 & 1 & \infty & \infty \\
 a & \infty & a & \infty \\
 \infty & \infty & \infty & \infty
\end{array}
\qquad
\begin{array}{c|ccc}
 \cdot & 1 & a & \infty \\ \hline
 1 & 1 & a & \infty \\
 a & a & \infty & \infty \\
 \infty & \infty & \infty & \infty .
\end{array}
\]
For $n\geq 2$, let $\M_n(S_7)$ be the semiring of $n\times n$ matrices
over $S_7$. Jiao and Ren proved that $\M_n(S)$ embeds into
$\M_{n+1}(S)$ for every additively idempotent semiring $S$ and every
$n\geq 2$. They asked whether the resulting chain for $S_7$ is constant
from dimension two onward and observed that their nilpotency result suggests
an affirmative answer \cite[Theorem~2.1 and Section~5]{JiaoRen}. We show
that the inclusion at $n=2$ is strict.

\begin{theorem}\label{thm:main}
The identity
\begin{equation}\label{eq:identity}
 x+y+x^2+y^2 \approx x+y+x^2+y^2+yx
\end{equation}
holds in $\M_2(S_7)$ and fails in $\M_n(S_7)$ for every $n\geq 3$.
Consequently,
\[
 \V(\M_2(S_7))\subsetneq\V(\M_n(S_7))
 \qquad(n\geq 3).
\]
In particular,
$\V(\M_2(S_7))\subsetneq\V(\M_3(S_7))$.
\end{theorem}

\begin{proof}
Let $X,Y\in\M_2(S_7)$ and set
\[
 P=X+Y+X^2+Y^2.
\]
We prove $P+YX=P$ entrywise. From the Cayley tables, a finite nonempty
sum in $S_7$ is $1$ precisely when all of its summands are $1$, and is
$a$ precisely when all of its summands are $a$. Also, a product is $1$
only when both factors are $1$, while $sa=as=a$ holds only for $s=1$.

Fix $i,j\in\{1,2\}$. If $P_{ij}=\infty$, then
$(P+YX)_{ij}=P_{ij}$. Suppose $P_{ij}=1$. Then
$(X^2)_{ij}=(Y^2)_{ij}=1$. Hence, for each $k\in\{1,2\}$,
\[
 Y_{ik}Y_{kj}=1
 \quad\text{and}\quad
 X_{ik}X_{kj}=1.
\]
Thus $Y_{ik}=X_{kj}=1$ for each $k$, so $(YX)_{ij}=1$.

Suppose instead that $P_{ij}=a$. Then
\[
 X_{ij}=Y_{ij}=(X^2)_{ij}=(Y^2)_{ij}=a.
\]
We have $i\neq j$, since otherwise the summand
$X_{ii}^2=a^2=\infty$ would force $(X^2)_{ii}=\infty$.
Since $i\neq j$ and the matrices are $2\times2$, the summation indices
are precisely $i$ and $j$, and therefore
\[
 a=(X^2)_{ij}=X_{ii}a+aX_{jj}.
\]
Both summands must be $a$, so $X_{ii}=X_{jj}=1$. The same argument gives
$Y_{ii}=Y_{jj}=1$. Consequently,
\[
 (YX)_{ij}=Y_{ii}X_{ij}+Y_{ij}X_{jj}=1\cdot a+a\cdot1=a.
\]
This proves \eqref{eq:identity} in $\M_2(S_7)$.

For failure in dimension three, take
\[
 X=\begin{pmatrix}
 1&1&1\\
 1&1&1\\
 a&a&1
 \end{pmatrix},
 \qquad
 Y=\begin{pmatrix}
 1&1&1\\
 a&1&1\\
 a&1&1
 \end{pmatrix}.
\]
At entry $(3,1)$,
\[
 X_{31}=Y_{31}=a,
 \qquad
 (X^2)_{31}=a\cdot1+a\cdot1+1\cdot a=a,
\]
\[
 (Y^2)_{31}=a\cdot1+1\cdot a+1\cdot a=a,
 \qquad
 (YX)_{31}=a\cdot1+1\cdot1+1\cdot a=\infty.
\]
Thus the two sides of \eqref{eq:identity} have respective $(3,1)$ entries
$a$ and $\infty$, so the identity fails in $\M_3(S_7)$.

By the embedding theorem of Jiao and Ren,
$\M_3(S_7)$ embeds into $\M_n(S_7)$ for every $n\geq 3$; hence
\eqref{eq:identity} fails in every such $\M_n(S_7)$. Their theorem also
gives $\M_2(S_7)\hookrightarrow\M_n(S_7)$, and therefore
$\V(\M_2(S_7))\subseteq\V(\M_n(S_7))$. Since an algebra and the variety
it generates satisfy the same identities, \eqref{eq:identity} makes the
inclusion strict.
\end{proof}

Theorem~\ref{thm:main} gives a negative answer to the question of Jiao and
Ren and disproves stabilization at $\V(\M_2(S_7))$. It does not distinguish
consecutive varieties from dimension three onward and therefore does not
determine whether the chain stabilizes at a later dimension.

\begin{thebibliography}{1}

\bibitem{JiaoRen}
J.~Jiao and M.~Ren,
\emph{The finite basis problem for matrix semirings $\M_n(S_7)$},
arXiv:2607.09677v1 [math.RA], 2026.

\end{thebibliography}

\end{document}